
As of the end of June 2025, 117.3 million people had been forced to flee their homes globally due to violence, conflict and human rights violation (UNHCR, 2025). When persecution is politically motivated, exile becomes politically selective. What are the political consequences for host communities when they receive ideologically motivated refugees? Contact between the displaced and the native population could not only reduce prejudice \citep{allport1954prejudice}, but also facilitate the diffusion of ideas between both groups. On the contrary, if the integration of the displaced population is limited, the ideas refugees bring may be perceived as a specific threat, generating political backlash. Perceptions of threat could be amplified when the displaced population is politically distant. Theoretically, the expected effect is therefore ambiguous.


To examine this question, we focus on the mass exodus from Spain to France during the final phase of the Spanish Civil War, known as La Retirada (``the retreat”). In early 1939, the collapse of Republican forces in Catalonia—who represented various left-wing groups including anarchists, communists, and other republican factions—triggered the flight of nearly half a million people. French authorities, unprepared for the scale of the inflow, treated the refugees as a security concern and responded by interning them. The locations were chosen under extreme time pressure based on logistical availability. Because no permanent facilities had been prepared in advance, the selection of camp sites was improvised. In the South, remote beaches, open plains, and abandoned barracks were selected primarily for their availability and their capacity to confine large numbers of people \citep{Coignard2012, dreyfus-armand_refugies_2016}. This process generated a geographically uneven yet plausibly exogenous distribution of camp locations.

Spanish refugees were regarded as ``undesirable foreigners” and seen as an economic burden. Therefore, propaganda and incentives were used to promote ``voluntary” repatriation to Spain or emigration to third countries. By early August, nearly half of those who had crossed the border during the episode had gone back \citep{perez_rodriguez_indeseables_2022}. Among those who remained, many were recruited into labour and military service. With the onset of World War II, their situation deteriorated further: some joined the French Resistance, while thousands were captured and deported to Nazi concentration camps. By the end of the World War II, 160,000 Republican exiles were still living in the country \citep{Coignard2012}. This allows us to study the effect of exposure to a refugee camp without the results being completely driven by subsequent changes in the demographic composition of the population or the economic structure, as most of the Spanish population eventually left the municipalities.

By leveraging this historical episode, we exploit the improvised placement of camps and estimate their effect on voting behaviour using a difference-in-differences strategy that accounts for spatial spillovers. Exposure to a refugee camp decreased the share of votes to the socialist party by 2,7 percentage points. Given that the refugees were themselves strongly aligned with left-wing and socialist ideals, this decline is consistent with a political backlash among host populations.

Yet the response was not uniform across the political spectrum. Several classic left-wing organizations—particularly labour and professional unions—were actively involved in assisting the refugees and thus maintained close contact with them. This heterogeneity is reflected in our findings. We observe a positive impact on support for the Communist Party - municipalities within 1 km away from the refugee camp increased their share of vote for the communist in 1,6 percentage points.In line with these results, municipalities that hosted refugees exhibit higher levels of mobilization against the Nazi occupation, indicating that exposure also fostered forms of political engagement aligned with the refugees’ ideological profile. Neighbouring municipalities show similar positive relationships, pointing to meaningful spillover effects beyond the immediate location of the camps.  


Our research relates to the broad literature on the political effect of immigration. Much of the existing literature emphasizes recent immigration trends and their immediate impact on public opinion and political behavior \citep{halla_immigration_2017, dustmann_refugee_2018, steinhardt_xenophobic_2018}. \citet{steinmayr_contact_2021} demonstrates that the type of contact matters: in Upper Austria, temporary refugee passage towards Germany in 2015 increased far-right voting, whereas longer-term settlement reduced it. In the French context, studies have found that larger immigrant populations tend to correlate with stronger electoral performance for the Front National at the département level \citep{edo_immigration_2019}. However, this pattern does not hold at the more localized municipal level, where the relationship appears to be reversed \citep{della2013competitive, vasilopoulos2022residential, vertier_dismantling_2023}. \citet{dazey_mosque_2024} finds that support for the Front National tends to rise in polling stations located at moderate distances from mosques, but declines beyond that range. We contribute by studying the consequences of an historical event, which allows us to explore the effect in the long-term.

Moreover, our study is also linked to the literature on migration and the transmission of political ideas \citep{spilimbergo_democracy_2009, barsbai_effect_2017, dippel_leadership_2021, calderon_racial_2023, bazzi_confederate_2023}. Previous research concentrates mostly on the effect at the origin country, with the exception of \citet{dippel_leadership_2021}, who study how the leaders of the failed German revolution of 1848–1849, expelled to the United States, positively influenced antislavery culture. We do not focus on the effects of a small political elite, as most government officials emigrated to Latin America \citep{alted2005vencidos}, but rather on a politically selected general population, which makes the analysis particularly relevant considering contemporary refugee flows. 